\documentclass[book.tex]{subfiles}

\begin{document}
\chapter{Generating Functions}
\label{chap:generating_functions}
\noindent\rule{11cm}{0.4pt} \\
\textbf{Prerequisites:} \autoref{chap:math_basics} \\
\textbf{Difficulty Level:} ** \\
\noindent\rule{11cm}{0.4pt}
\vspace{5mm}

Generating functions are tools for converting sequences into functions. There are a few choices for performing this conversion. Ordinary generating functions use an expansion in powers of $x$

\begin{align}
    G(a_n; x) &= \sum_{n=0}^\infty a_n x^n
\end{align}

\section{Generating Functions in Probability}

To ease calculations, we define the \textit{generating function} of a probability distribution as 
\begin{align}
    g(k) = \sum_{n=0}^{\infty}e^{-nk}P_n
\end{align}
which gives:

\begin{align}\label{Eq14}
\langle n^m \rangle = \lim_{k\rightarrow 0}\left[(-1)^m\frac{d^mg(k)}{dk^m}\right]
\end{align}

\noindent which alters the calculation to the evaluation of derivatives.

For the Poisson distribution, we have:

\begin{equation}\label{Eq15}
g(k)= \sum_{n=0}^{\infty}e^{-kn}\frac{a^n}{n!}e^{-a}=exp\left[a\left(e^{-k}-1\right)\right]
\end{equation}

Using the generating function, we have:

\begin{align}
&\langle n \rangle = -\lim_{k \rightarrow 0} \left\{-ae^{-k} exp \left[a \left(e^{-k}-1\right)\right]\right\}=a \notag\\
&\langle n^2 \rangle = -\lim_{k\rightarrow 0} \left\{ \left(a e^{-k}+a^2 e^{-2k}\right) exp\left[ a \left(e^{-k}-1\right)\right] \right\} = a + a^2 \notag
\end{align}

%\section{Exercises}
%\begin{enumerate}
%    \item TODO
%\end{enumerate}

\end{document}